% Seccion: Muestreo y aliasing
% Responsable: R1. Extension: 1,0 pagina(s).
% Solo R1 edita este archivo.
%
% Que debe contener (segun el enunciado y la pauta):
%   - Deducción y justificación de la expresión de frecuencia_alias.
%   - Figura muestreo_alias.png y tabla muestreo_alias.csv: frecuencia medida vs. predicha.

\section{Muestreo y aliasing}
\label{sec:muestreo}

\paragraph{Frecuencia aparente.}
Una senoidal de frecuencia $f$ muestreada a $f_s$ da $x[n]=\cos(2\pi f n/f_s)$. Las muestras no cambian
(i)~al sumar a $f$ un múltiplo de $f_s$, porque $\cos(2\pi (f+kf_s)n/f_s)=\cos(2\pi fn/f_s+2\pi kn)$ y $kn$
es entero, ni (ii)~al cambiar su signo, porque el coseno es par. Todas las frecuencias $\pm f+kf_s$
producen entonces exactamente las mismas muestras y no se pueden distinguir. Con $k^\ast=\operatorname{round}(f/f_s)$,
la diferencia $f-k^\ast f_s$ cae en $[-f_s/2,\,f_s/2]$ y, por paridad, el representante del conjunto en
$[0,\,f_s/2]$ es
\begin{equation}
  f_a \;=\; \bigl|\,f - f_s\,\operatorname{round}(f/f_s)\,\bigr| ,
  \label{eq:alias}
\end{equation}
que es lo que calcula \texttt{frecuencia\_alias}. Si $f<f_s/2$ entonces $k^\ast=0$ y $f_a=f$ (sin aliasing,
criterio de Nyquist); si no, la frecuencia se refleja alrededor de los múltiplos de $f_s/2$.

\paragraph{Verificación.}
Se muestrearon senoidales de $f\in\{3,7,13,23\}$~Hz con $f_s=10$~Hz durante 10~s y se midió la frecuencia
aparente como el máximo del espectro (\texttt{espectro}, ventana rectangular, resolución
$\Delta f = 1/T = 0{,}1$~Hz). La ecuación~\eqref{eq:alias} predice 3~Hz en los cuatro casos, ya que
$7=10-3$, $13=10+3$ y $23=2\cdot10+3$; la medición coincide (tabla~\ref{tab:muestreo}). El error exacto
de cero se debe a que cada tono cae justo en un bin ($f$ y $T$ enteros); en general el error de la medición
está acotado por la resolución $\Delta f$.

% Valores de out/tablas/muestreo_alias.csv (generada por figuras/fig_muestreo.py).
\begin{table}[H]
  \centering
  \caption{Frecuencia aparente de senoidales muestreadas a $f_s=10$~Hz: predicha por
           \eqref{eq:alias} y medida con el espectro.}
  \label{tab:muestreo}
  \begin{tabular}{rrrr}
    \toprule
    $f$ [Hz] & $f_a$ predicha [Hz] & $f_a$ medida [Hz] & error [Hz] \\
    \midrule
    3  & 3,0 & 3,0 & 0,0 \\
    7  & 3,0 & 3,0 & 0,0 \\
    13 & 3,0 & 3,0 & 0,0 \\
    23 & 3,0 & 3,0 & 0,0 \\
    \bottomrule
  \end{tabular}
\end{table}

\figura{muestreo_alias.png}{0.72}{Las senoidales de 3~Hz y 13~Hz (curvas, muestreadas muy fino) pasan por
  exactamente las mismas muestras a $f_s=10$~Hz, pues $13=3+f_s$.}{fig:muestreo}

La figura~\ref{fig:muestreo} lo muestra para 3~Hz y 13~Hz: a $f_s=10$~Hz las muestras son idénticas, y por lo
tanto ningún procesamiento posterior puede recuperar cuál de las dos señales había. Las pruebas verifican
propiedades, no valores: $f_a\in[0,f_s/2]$, periodicidad en $f_s$, paridad e idempotencia de
\eqref{eq:alias}; que $f$ y $f_a$ generen las mismas muestras; y que, para 30 pares $(f,f_s)$ aleatorios,
el pico del espectro caiga en \texttt{frecuencia\_alias}.
